\chapter{The room and the airflow simulations}
\label{ch:room}

\section{The classroom}
The classroom is 15.53~m long, 9.16~m wide and 3.15~m high. The geometry comes from STL files of the
real room. One long wall has 8 windows (W1--W8, each 0.61~m wide and as high as the room in the model),
the opposite wall has 2 doors (1.06~m wide, 2.11~m high), and 4 round columns stand in the room
(\fref{fig:room}). Fresh air comes in through the open windows and leaves through the open doors.

\begin{figure}[h]
  \centering
  \figplaceholder{top view of the room with windows, doors, columns and seating area}
  \caption{The classroom. \todo{draw}}
  \label{fig:room}
\end{figure}

The people sit in a seating area of $9.5\times5$~m between 1.0 and 1.2~m height, where the \cotwo{} is
released. This area is still a placeholder until the real seating plan is known
\todo{update when the seating plan is available}. All simulations use 20 people; because
\eref{eq:co2} is linear, the results can be scaled to any number of people.

\section{Window settings}
In the first setup, the 8 windows were grouped into three zones (W1--W3, W4--W6, W7--W8), and both doors were
open. I simulated 39 settings with this setup (31 for training, 8 for testing).

Later we changed the setup to the way the room will be used with the sensors: only windows W1, W3 and W6 can
open, door~1 is open and door~2 is closed. Each of the three windows is closed or open with an inflow speed
between 0.2 and 2.0~m/s. I chose 28 settings (\tref{tab:settings}): each window alone at 0.5, 1 and 2~m/s,
all three together at the same speeds, and 16 random combinations. 8 of the random ones are kept for
testing and never used in training. This thesis mainly reports the current setup.

\begin{table}[h]
  \centering
  \caption{The 28 window settings of the current setup.}
  \label{tab:settings}
  \begin{tabular}{l>{\raggedright}p{7.6cm}l}
    \toprule
    Settings & Windows & Use\\
    \midrule
    N01--N09 & one window open (W1, W3 or W6) at 0.5, 1 or 2~m/s & training\\
    N10--N12 & all three open at 0.5, 1 or 2~m/s & training\\
    N13--N28 & random speeds 0.2--2.0~m/s, each window closed with probability 0.3 & 8 training, 8 test\\
    \bottomrule
  \end{tabular}
\end{table}
The test settings are N14, N15, N19, N20, N23, N24, N27 and N28.

\section{Simulation setup}
\tref{tab:sim} summarises the OpenFOAM settings. Each open window is an inlet with a uniform velocity normal
to the wall; the speed rises smoothly from zero and reaches the full value after about 2~s
($v(t)=v_0\tanh(1.5\,t/\mathrm{s})$). The
inlet turbulence intensity is 5\,\% with a length scale of 0.1~m. The doors are pressure outlets ($p=0$) that
also allow backflow. Walls, columns, closed windows and the closed door are no-slip walls with wall functions.
The \cotwo{} is also computed by OpenFOAM itself (scalar transport with $\mathrm{Sc}_t=0.7$), which I use as
the reference for my own solver.

\begin{table}[h]
  \centering
  \caption{OpenFOAM settings.}
  \label{tab:sim}
  \begin{tabular}{>{\raggedright}p{4.2cm}p{9.5cm}}
    \toprule
    Item & Setting\\
    \midrule
    Solver & OpenFOAM v2412, \texttt{pimpleFoam} (transient), k-$\omega$ SST with wall functions\\
    Mesh & hexahedral, 8~cm cells, 856,284 cells after removing the columns\\
    Air & kinematic viscosity $1.5\times10^{-5}$~m\textsuperscript{2}/s\\
    Time step & adaptive, Courant number $\le 1$\\
    Simulated time & 600~s\\
    \cotwo{} source & 20 people, 6~L/min at 40,000~ppm each, spread evenly over the seating area\\
    \cotwo{} diffusivity & $D=1.6\times10^{-5}$~m\textsuperscript{2}/s\\
    \bottomrule
  \end{tabular}
\end{table}

\section{Simplifications of the geometry}
The geometry of the model is a simplified version of the STL files of the real room. Compared with
the STL files, the model differs in these points:
\begin{itemize}
  \item a solid block on the door wall between the two doors ($x=7.39$--$8.23$~m, 0.72~m deep,
        floor to ceiling, probably a shaft) is missing. It lies about 5~m from each door, where the
        air moves slowly, so I expect its effect to be local, but I have not measured it;
  \item the windows sit in recesses of 0.15--0.3~m in the real room; in the model they are flat
        openings in a flat wall;
  \item the walls of the real room are slightly stepped, and the ceiling is at 3.13~m instead of
        3.15~m;
  \item every window is an opening of the full room height (0.61~m $\times$ 3.15~m). The real
        windows are smaller, so a measured window speed has to be converted through the volume flow
        before it is used as model input (\cref{ch:sensors});
  \item when a door is closed, the gap under it (about 1.5~m $\times$ 0.4~cm, estimated by the
        supervisor) is neglected;
  \item the air has a constant temperature: heat from people and walls and the resulting
        buoyancy are not modelled.
\end{itemize}
Adding the block or the recesses would mean running all OpenFOAM cases again, so I keep them as
simplifications.

\section{Mesh study}
The first setup used a 10~cm mesh. To check whether this is fine enough, I simulated two settings on meshes
with 15, 10 and 8~cm cells and compared them with the 8~cm mesh (\tref{tab:mesh}). For the faster setting
(S08) the 10~cm mesh gives about 30\,\% less \cotwo{} in the room and a less steady flow. The current setup
therefore uses the 8~cm mesh.

\begin{table}[h]
  \centering
  \caption{Mesh study: difference to the 8~cm mesh at 600~s (relative $\ell_2$ difference of the
  velocity in the room and of the \cotwo{} at 1.6~m, once above the outdoor level and once in the
  paper's metric with the 400~ppm included; total \cotwo{} in the room relative to the 8~cm mesh).}
  \label{tab:mesh}
  \begin{tabular}{llrrrrr}
    \toprule
    Setting & Mesh & Cells & Velocity & \cotwo{} at 1.6~m & Paper metric & Total \cotwo{}\\
    \midrule
    S02 & 15~cm & 132,258 & 25.8\,\% & 47.3\,\% & 5.47\,\% & 1.004\\
        & 10~cm & 438,929 & 22.7\,\% & 42.4\,\% & 4.91\,\% & 0.931\\
    S08 & 15~cm & 132,258 & 28.6\,\% & 36.7\,\% & 1.54\,\% & 0.749\\
        & 10~cm & 438,929 & 40.4\,\% & 38.3\,\% & 1.60\,\% & 0.697\\
    \bottomrule
  \end{tabular}
\end{table}

The wall resolution ($y^+$ of the first cell, 8~cm mesh) is between 2 and 140 for S02 and between 21
and 571 for S08, with averages of 42 and 163. The wall functions of OpenFOAM v2412 cover both ranges:
the wall function for $\omega$ blends its viscous and its logarithmic value, and the wall function for
the turbulent viscosity sets it to zero below $y^+\approx11$ (viscous sublayer) and uses the log law
above. They are least accurate in the buffer layer ($5\lesssim y^+\lesssim30$). Most wall cells are in
the log-law range, and the flow directly at the walls is not the quantity of interest in this study.
\todo{mention checkMesh result for the 8 cm mesh}
